\documentclass[12pt]{article}
\usepackage{inputenc}
\usepackage[spanish]{babel}
\decimalpoint
\usepackage{geometry}
\usepackage{amsmath}
\usepackage{amssymb}
\usepackage{amsthm}
\usepackage{hyperref}
\usepackage{tikz}
\usepackage{titling}

\newgeometry{left=0.6in, right=0.6in, top=0.4in, bottom=0.4in}

\newtheorem{problem}{Problema}

\title{Procesos Estocásticos I | Examen Parcial 3}
\author{Semestre 2026-2}
\date{25 de mayo de 2026}

\begin{document}

\maketitle
\vspace{-0.8cm}
\small
\textbf{Instrucciones:} Resuelva los siguientes problemas. Encierre sus respuestas finales en un recuadro. Justifique cada paso de sus cálculos y razonamientos.

% ---------------------------------------------------------------
% \begin{problem}
% Se propone el siguiente sistema de apuestas. Un jugador lanza repetidamente un dado
% equilibrado con caras $\{1,2,3,4,5,6\}$; sea $D_n$ el resultado de la $n$-ésima tirada y
% $\mathcal{F}_n = \sigma(D_1,\ldots,D_n)$ la filtración natural. La fortuna acumulada del jugador
% sigue la recurrencia
% \[
%     M_0 = 0, \qquad M_n = M_{n-1} + H_n\,\xi_n, \quad n \geq 1,
% \]
% donde
% \[
%     H_n = \mathbf{1}\!\bigl[D_{n+1} \leq 3\bigr], \qquad
%     \xi_n =
%     \begin{cases}
%         +2 & \text{si } D_n \in \{1,2\}, \\[2pt]
%         +1 & \text{si } D_n \in \{3,4\}, \\[2pt]
%         -4 & \text{si } D_n \in \{5,6\}.
%     \end{cases}
% \]
% \begin{enumerate}
%     \item[(a)] Determine si $H_n$ es $\mathcal{F}_{n-1}$-medible. ¿Qué condición de la definición
%         de martingala viola este hecho?
%     \item[(b)] Calcule $\mathbb{E}[\xi_n]$. ¿Favorece el juego al jugador o a la casa?
%     \item[(c)] Sea $\mathcal{G}_n = \sigma(D_1,\ldots,D_{n+1})$. Con respecto a $\{\mathcal{G}_n\}$,
%         el proceso $\{M_n\}$ sí es adaptado. Calcule $\mathbb{E}[M_n \mid \mathcal{G}_{n-1}]$ y
%         determine si $\{M_n,\mathcal{G}_n\}$ es martingala, supermartingala o submartingala.
% \end{enumerate}
% \end{problem}

% ---------------------------------------------------------------
\begin{problem}
Sea $\{N_t\}_{t \geq 0}$ un proceso de Poisson de tasa $\lambda > 0$ definido sobre la filtración natural del proceso
$\{\mathcal{F}_t\}$.
\begin{enumerate}
    \item[(a)] Demuestre que el proceso compensado $\{\widetilde{N}_t\}_{t\geq 0}$ dado por $\widetilde{N}_t := N_t - \lambda t$ es una
        martingala con respecto a $\{\mathcal{F}_t\}$.
    \item[(b)] Sea $W_n = \inf\{t \geq 0 : N_t = n\}$ el instante de la $n$-ésima llegada.
        Verifique que $W_n$ es un tiempo de paro y que $\{\widetilde{N}_{t\wedge W_n}\}$
        satisface las hipótesis del Teorema del Paro Opcional. Aplique dicho teorema para
        calcular $\mathbb{E}[W_n]$ y corrobore con lo conocido sobre la distribución de $W_n$.
\end{enumerate}
\end{problem}

% ---------------------------------------------------------------
\begin{problem}
Sean $\{N_t^{(1)}\}_{t \geq 0}$ y $\{N_t^{(2)}\}_{t \geq 0}$ procesos de Poisson independientes con tasas $\lambda_1 > 0$ y
$\lambda_2 > 0$, respectivamente. Sea $N_t = N_t^{(1)} + N_t^{(2)}$.
\begin{enumerate}
    \item[(a)] Calcule $\mathbb{P}(N_t^{(1)} = k \mid N_t = n)$ para $k = 0, 1, \ldots, n$ e
        identifique la distribución resultante indicando sus parámetros.
\end{enumerate}
\end{problem}

% ---------------------------------------------------------------
\begin{problem}
Las solicitudes a un servidor de cómputo llegan de acuerdo con un proceso de Poisson $\{N_t\}_{t \geq 0}$ de tasa
$\lambda = 3$ solicitudes por minuto. Sea $N_t$ el número de solicitudes recibidas en $[0,t]$.
\begin{enumerate}
    \item[(a)] Calcule $\mathbb{P}\!\left(N_4-N_2>5 \;\Big|\; N_1=2\right)$. Exprese su
        respuesta en términos de la función de distribución acumulada de una distribución Poisson
        con el parámetro que corresponda.
    \item[(b)] Dado que se recibieron exactamente $3$ solicitudes en el intervalo $[0,2]$,
        determine la distribución del tiempo de llegada de la primer solicitud. Identifique la
        distribución e indique sus parámetros. Hint: Use la regla de Bayes sobre la definición de densidad condicional de $T_1$ dado $N_2=3$.
\end{enumerate}
\end{problem}

% ---------------------------------------------------------------
%\begin{problem}
%Sean $\{B_t^{(1)}\}$ y $\{B_t^{(2)}\}$ movimientos brownianos estándar independientes.
%\begin{enumerate}
%    \item[(a)] Demuestre que el proceso
%        \[
%            W_t = \frac{B_t^{(1)} + B_t^{(2)}}{\sqrt{2}}
%        \]
%        es un movimiento browniano estándar.
%    \item[(b)] El logaritmo del rendimiento anual de una acción se modela como $X = \sigma B_1$,
%        con $\sigma = 0.20$ y $B_1 \sim \mathcal{N}(0,1)$. Encuentre $L$ y $U$ tales que
%        $\mathbb{P}(L \leq X \leq U) = 0.95$, expresando $L$ y $U$ en términos de $\sigma$ y del
%        cuantil $z_{0.025} \approx -1.96$. Interprete el intervalo $[L,U]$ en el contexto
%        financiero.
%\end{enumerate}
%\end{problem}

% ---------------------------------------------------------------
\begin{problem}
Sea $\{B_t\}_{t \geq 0}$ un movimiento browniano estándar.
\begin{enumerate}
    \item[(a)] Demuestre que $\mathbb{P}(B_t \geq B_s) = \tfrac{1}{2}$ para cualesquiera tiempos
        $0 \leq s < t$.
    \item[(b)] Sea $s > 0$ fijo y defina $X_t = B_s - B_{s-t}$ para $t \in [0,s]$. Demuestre que
        $\{X_t : t \in [0,s]\}$ es un movimiento browniano estándar en $[0,s]$.
 \end{enumerate}
\end{problem}

\end{document}
